\chapter*{Abstract}
\addcontentsline{toc}{chapter}{Abstract}

Quantum information processing depends on the preparation and preservation of coherent quantum states, yet realistic systems are inevitably subject to environmental noise. This dissertation-style study analyzes the \emph{Quantum Noise Sensitivity Mapping Engine}, a compact Python simulator designed to study state degradation using density matrices and explicitly implemented Kraus channels. The engine models a single-qubit superposition state, a Bell state, and GHZ states with two to four qubits, and it applies local amplitude damping and phase damping noise without relying on external quantum-computing frameworks. The simulation pipeline constructs ideal density matrices, embeds single-qubit Kraus operators into multipartite Hilbert spaces, applies noise independently to each qubit, validates the resulting density matrices, and computes fidelity and purity. Parameter sweeps over noise strength and a GHZ scaling study over qubit number are exported as CSV files and SVG figures. The numerical outputs show that fidelity decreases as noise increases in all bundled studies, while purity distinguishes the physical action of the two channels: phase damping drives states toward mixedness, whereas strong amplitude damping can restore purity by relaxing the system toward a pure ground state even as fidelity to the target state continues to decline. The GHZ study additionally shows a consistent decrease in robustness with increasing qubit number under phase damping. The project demonstrates that a small, transparent, framework-free implementation can still support a technically meaningful study of open-system quantum behavior while remaining readable, extensible, and pedagogically valuable.
